\begin{formbox}
\textbf{Name:} Jingyan Wei \hfill \textbf{Student ID:} 2361411\\
\textbf{Role this project:} Mathematical theory --- variational formulation
and the FDM/FEM comparison (Role 2).

\textbf{The specific things I did this project:}
\begin{enumerate}[leftmargin=2em,itemsep=0.25em,topsep=0.15em]
    \item I worked out the quadrature behind the P1 stiffness and load
    integrals in Section~3.3, \path{project2/report/algorithm_section.tex}.
    Each square is split along its diagonal, so every element integral of
    $\int_T a\,\nabla\varphi_i\cdot\nabla\varphi_j\,dx$ is evaluated on the
    reference triangle after a Duffy map, with barycentric coordinates
    $((1-r)(1-s),r,(1-r)s)$ and Jacobian factor $1-r$. I checked the degree
    that the $5\times5$ Gauss--Legendre rule must integrate: the reference
    weights sum to $1/2$, and for an affine coefficient $a$ --- and for the
    interface-aligned piecewise-constant $a$ --- the stiffness entries come
    out exact to round-off. A higher-order rule independently checks the
    smooth load and the continuous $L^2$ error integral on a representative
    mesh.

    \item I defined how a nodal FDM solution is put into the same norm as the
    P1 solution, which the depth comparison needs because the two
    discretisations do not even live in the same function space. Both nodal
    solutions are reconstructed by P1 interpolation on the common triangles
    and integrated against the exact solution with a $6\times6$ Duffy rule
    on the mesh where the comparison is reported. I also recorded the parity
    bookkeeping of the layered experiment: FDM uses $N=31,63,127,255$ while
    the interface-fitted FEM uses $N=32,64,128,256$, so the two structural
    checks are read on their own meshes instead of being forced into one
    refinement table.

    \item I analysed the conditioning and iteration tables
    (the condition-number table of Section~4.4). At $N=63$ the
    measured $\kappa_2(A_h)$ rises from $1.61\times10^3$ to
    $7.15\times10^3$ to $6.49\times10^4$ as the contrast goes from $1$ to
    $10$ to $100$, while unpreconditioned CG iterations rise from $216$ to
    $521$ to $1354$ but Jacobi--CG stays at $216$, $224$, $226$. I checked
    that the growth with mesh size is consistent with the $O(h^{-2})$ upper
    bound, and I kept the writing to what the numbers support: they justify
    the solver choice and the use of Jacobi scaling on this layered geometry,
    and they do not claim contrast-independent performance for arbitrary
    heterogeneous media.
\end{enumerate}

\textbf{What I would do differently:} I would fix the quadrature degree
that each integral needs by a patch test before writing the integral down,
rather than checking it afterwards, and I would compute the largest
eigenvalue by an independent route before a table was built on a single
estimate --- a spectral estimate that looks unusually clean is a reason to
check it, not a reason to trust it.

\textbf{Anything I should have asked for help with earlier:} I should have
asked at the first verification clinic exactly which norm each table
reports, nodal or reconstructed, because I initially compared two checks
that were written in different norms and the agreement I saw was not
meaningful. I also should have asked for the assembler's return convention
(vertex ordering, boundary set) before reading for it in the code.

\textbf{AI use:} the AI-assisted work that reached the report is logged in
Part~A of the AI Transparency Log (audits~1--3). The numbers quoted above
are the ones reproduced by the committed tables, not the ones in a chat
transcript.

\vfill
\textbf{Signature:} \underline{\hspace{4.5cm}} \hfill
\textbf{Date:} \underline{\hspace{2.5cm}}
\end{formbox}
